% Noi dung Chuong 1 (tam thoi bo khoi main.tex ngay 19/09/2026).
% Muon khoi phuc: dan lai toan bo phan duoi vao ngay sau \label{sec:overview} trong main.tex.

\subsection{Nguồn gốc và quy mô}
Nhóm được cung cấp một bản sao lưu cơ sở dữ liệu tên \texttt{CompanyX.bak}, dung lượng 214 MB. Khi khôi phục bản sao lưu này lên SQL Server, tên gốc bên trong hiện ra là \texttt{AdventureWorks2022}. CompanyX vì vậy chính là bộ dữ liệu mẫu AdventureWorks của Microsoft, mô phỏng công ty Adventure Works Cycles: một doanh nghiệp sản xuất và bán xe đạp, linh kiện, quần áo và phụ kiện đạp xe. Công ty bán hàng theo hai kênh, trực tuyến cho khách cá nhân và bán sỉ qua các cửa hàng đại lý do nhân viên kinh doanh phụ trách, trên ba khu vực Bắc Mỹ, Châu Âu và Châu Úc. Bảng~\ref{tab:overview} tóm tắt quy mô kho dữ liệu.

\begin{table}[H]
\centering
\caption{Quy mô kho dữ liệu CompanyX}
\label{tab:overview}
\begin{tabularx}{\textwidth}{>{\raggedright\arraybackslash}p{5.5cm} X}
\toprule
\rowcolor{tblhead} \textbf{Thông số} & \textbf{Giá trị} \\
\midrule
Số bảng dữ liệu & 72 bảng, chia thành 6 nhóm nghiệp vụ \\
Số view (truy vấn được đặt tên, không lưu dữ liệu riêng) & 24 \\
Khoảng thời gian có giao dịch & 31/05/2011 -- 30/06/2014 (37 tháng) \\
Số khách hàng có mua hàng & 19.119 \\
Số đơn hàng & 31.465 đơn, gồm 121.317 dòng sản phẩm \\
Tổng doanh thu & 123,16 triệu USD \\
Số sản phẩm & 504 sản phẩm, thuộc 4 danh mục lớn và 37 danh mục con \\
Số lãnh thổ bán hàng & 10 \\
\bottomrule
\end{tabularx}
\end{table}

\subsection{Sáu nhóm bảng theo bộ phận nghiệp vụ}
72 bảng được xếp thành 6 nhóm, mỗi nhóm ứng với một bộ phận trong công ty (Bảng~\ref{tab:schemas}). Với chủ đề khách hàng rời bỏ, hai nhóm Bán hàng (Sales) và Cá nhân (Person) chứa gần như toàn bộ thông tin cần thiết: khách là ai và đã mua gì. Nhóm Sản xuất (Production) cho biết tên và loại sản phẩm. Ba nhóm còn lại ghi nhận việc mua nguyên liệu từ nhà cung cấp, nhân sự và các bảng hệ thống, hầu như không liên quan đến chủ đề.

\begin{table}[H]
\centering
\caption{Các nhóm bảng trong CompanyX và mức độ liên quan đến chủ đề}
\label{tab:schemas}
\small
\begin{tabularx}{\textwidth}{l c X c}
\toprule
\rowcolor{tblhead} \textbf{Nhóm} & \textbf{Số bảng} & \textbf{Chứa thông tin gì} & \textbf{Liên quan} \\
\midrule
Sales (Bán hàng) & 19 & Khách hàng, cửa hàng đại lý, đơn hàng, chi tiết đơn, lãnh thổ, nhân viên kinh doanh, chương trình khuyến mãi, lý do mua, thẻ tín dụng, tỷ giá & Cao \\
Person (Cá nhân) & 13 & Thông tin cá nhân, địa chỉ, email, điện thoại, quốc gia/tỉnh thành & Cao \\
Production (Sản xuất) & 25 & Sản phẩm, danh mục, mẫu xe, tồn kho, lệnh sản xuất, lịch sử giá & Trung bình \\
Purchasing (Mua hàng) & 5 & Nhà cung cấp, đơn mua nguyên liệu, phương thức vận chuyển & Thấp \\
HumanResources (Nhân sự) & 6 & Nhân viên, phòng ban, lương, ca làm việc & Thấp \\
dbo (Hệ thống) & 4 & Nhật ký hệ thống, phiên bản; kèm 4 view bổ sung riêng của CompanyX & Hỗ trợ \\
\bottomrule
\end{tabularx}
\end{table}

\subsection{Các bảng quan trọng và ý nghĩa của chúng}
Bảng~\ref{tab:tables} liệt kê những bảng được dùng trong phân tích, số dòng thực tế và thông tin mà mỗi bảng mang lại.

\begin{table}[H]
\centering
\caption{Các bảng chính phục vụ phân tích khách hàng}
\label{tab:tables}
\small
\begin{tabularx}{\textwidth}{>{\raggedright\arraybackslash}p{4.3cm} r X}
\toprule
\rowcolor{tblhead} \textbf{Bảng} & \textbf{Số dòng} & \textbf{Cho biết điều gì} \\
\midrule
Sales.Customer & 19.820 & Danh sách khách hàng. Mỗi khách hoặc là một cá nhân, hoặc là một cửa hàng đại lý, và thuộc về một lãnh thổ bán hàng \\
Sales.Store & 701 & Hồ sơ đại lý: tên cửa hàng, nhân viên kinh doanh phụ trách, kết quả khảo sát (loại hình, chuyên môn, doanh số năm, số nhân viên, năm mở cửa) \\
Person.Person & 19.972 & Hồ sơ cá nhân: họ tên, vai trò (khách hàng, nhân viên, nhà cung cấp\dots), có đồng ý nhận email khuyến mãi hay không, và kết quả khảo sát nhân khẩu học \\
Sales.SalesOrderHeader & 31.465 & Mỗi đơn hàng: ngày đặt, ngày giao, trạng thái (đã giao/đã hủy), kênh bán (online hay đại lý), nhân viên bán, tổng tiền \\
Sales.SalesOrderDetail & 121.317 & Từng sản phẩm trong đơn: sản phẩm gì, số lượng, đơn giá, mức giảm giá, chương trình khuyến mãi áp dụng \\
Sales.SalesOrderHeader\-SalesReason & 27.647 & Lý do khách chọn mua (giá, khuyến mãi, chất lượng, quảng cáo TV\dots) gắn với từng đơn \\
Sales.SalesTerritory & 10 & 10 lãnh thổ bán hàng, thuộc 3 khu vực: Bắc Mỹ, Châu Âu, Thái Bình Dương \\
Sales.SpecialOffer & 16 & Các chương trình khuyến mãi: giảm theo số lượng, theo mùa, sản phẩm mới, xả tồn kho \\
Production.Product & 504 & Sản phẩm và giá niêm yết; được xếp vào 37 danh mục con và 4 danh mục lớn (Xe đạp, Linh kiện, Quần áo, Phụ kiện) \\
Person.Address, CountryRegion & 19.614 / 238 & Địa chỉ và quốc gia của khách hàng \\
\bottomrule
\end{tabularx}
\end{table}

Kết quả khảo sát khách cá nhân (tuổi, thu nhập, nghề nghiệp\dots) không được lưu thành các cột riêng mà gói chung trong một trường XML của bảng \texttt{Person.Person}. CompanyX cung cấp sẵn view \texttt{Sales.vPersonDemographics} để tách trường này thành 12 thuộc tính: ngày sinh, giới tính, tình trạng hôn nhân, thu nhập năm, học vấn, nghề nghiệp, tổng số con, số con còn ở nhà, sở hữu nhà, số xe hơi, ngày mua đầu tiên và tổng mua trong năm. Khảo sát cửa hàng được lưu và tách tương tự trong view \texttt{Sales.vStoreWithDemographics}. Cần phân biệt: view là một truy vấn được đặt tên, đọc dữ liệu từ bảng gốc mỗi khi được gọi, không phải bảng lưu sẵn kết quả; vì vậy số liệu lấy qua view luôn khớp với bảng gốc, nhưng chi phí truy vấn nằm ở bước tách XML.

So với bộ AdventureWorks gốc, CompanyX có thêm bốn view trong nhóm \texttt{dbo}: \texttt{PersonCustomerView} (khách cá nhân kèm họ tên), \texttt{StoreCustomerView} (khách đại lý kèm tên cửa hàng), \texttt{SalesPersonView} (hồ sơ nhân viên kinh doanh với chỉ tiêu, thưởng, hoa hồng, doanh số) và \texttt{SalesOrderDetailView} (chi tiết đơn hàng kèm lý do mua, chỉ lấy những dòng có áp dụng khuyến mãi). Hai view đầu tách sẵn hai nhóm khách hàng, view cuối tập trung vào các dòng hàng có khuyến mãi.

\subsection{Những bảng nào thực sự cần cho chủ đề?}
Khi khảo sát sơ bộ, nhóm nhận thấy hai cụm bảng liên quan trực tiếp đến chủ đề: cụm Bán hàng gồm \texttt{Customer}, \texttt{SalesOrderHeader}, \texttt{SalesOrderDetail}, \texttt{SalesTerritory}, \texttt{Store}, và cụm Cá nhân gồm \texttt{Person}. Để chắc chắn, nhóm kiểm tra lại các mối liên kết được khai báo trong cơ sở dữ liệu (Bảng~\ref{tab:fk}). Kết quả xác nhận nhận định ban đầu: các bảng này nối với nhau thành một chuỗi liền mạch, đi từ khách hàng là ai đến họ đã mua gì.

\begin{table}[H]
\centering
\caption{Các mối liên kết giữa những bảng trung tâm}
\label{tab:fk}
\footnotesize
\begin{tabularx}{\textwidth}{>{\raggedright\arraybackslash}p{3.8cm} >{\raggedright\arraybackslash}p{2.4cm} >{\raggedright\arraybackslash}p{4.1cm} X}
\toprule
\rowcolor{tblhead} \textbf{Từ bảng} & \textbf{Qua cột} & \textbf{Đến bảng} & \textbf{Ý nghĩa} \\
\midrule
Sales.Customer & PersonID & Person.Person & Khách hàng là một cá nhân \\
Sales.Customer & StoreID & Sales.Store & Khách hàng là một đại lý \\
Sales.Customer & TerritoryID & Sales.SalesTerritory & Khách thuộc lãnh thổ nào \\
Sales.SalesOrderHeader & CustomerID & Sales.Customer & Đơn hàng của khách nào \\
Sales.SalesOrderHeader & TerritoryID & Sales.SalesTerritory & Đơn bán ở lãnh thổ nào \\
Sales.SalesOrderHeader & SalesPersonID & Sales.SalesPerson & Nhân viên nào bán (kênh đại lý) \\
Sales.SalesOrderDetail & SalesOrderID & Sales.SalesOrderHeader & Dòng sản phẩm thuộc đơn nào \\
Sales.SalesOrderDetail & ProductID, SpecialOfferID & Sales.SpecialOffer\-Product & Sản phẩm nào, khuyến mãi nào \\
Sales.SalesOrderHeader\-SalesReason & SalesOrderID, SalesReasonID & Sales.SalesOrderHeader, Sales.Sales\-Reason & Lý do mua của đơn \\
Sales.Store & SalesPersonID & Sales.SalesPerson & Đại lý do ai phụ trách \\
\bottomrule
\end{tabularx}
\end{table}

Về bản số (cardinality): một khách có nhiều đơn, một đơn có nhiều dòng hàng (quan hệ một--nhiều). Riêng lý do mua là quan hệ nhiều--nhiều: bảng cầu \texttt{SalesOrderHeaderSalesReason} có khóa chính ghép (\texttt{SalesOrderID}, \texttt{SalesReasonID}), nên một đơn có thể gắn nhiều lý do và mỗi lý do gắn với nhiều đơn. Khi nối bảng đơn hàng với bảng này, số dòng sẽ tăng theo số lý do của đơn; vì vậy khi thống kê theo lý do mua cần đếm riêng số lý do, số đơn và số khách để một đơn không bị tính nhiều lần.

Để trả lời trọn vẹn các câu hỏi về khách hàng, cần thêm ba nguồn ngoài danh sách trên: bảng lý do mua để biết vì sao khách mua; chuỗi \texttt{Product} $\to$ \texttt{ProductSubcategory} $\to$ \texttt{ProductCategory} để biết khách mua nhóm sản phẩm gì; và hai view nhân khẩu học đã nêu. Tất cả gộp lại khoảng 10 bảng, được vẽ lại theo ngôn ngữ nghiệp vụ trong Hình~\ref{fig:erd}. Bảng Khách hàng nằm ở giữa; mỗi khách có nhiều Đơn hàng; mỗi đơn gồm nhiều Dòng sản phẩm. Khách cá nhân được mô tả thêm bằng Hồ sơ cá nhân, khách đại lý bằng Hồ sơ cửa hàng, và mỗi đơn hàng gắn với một Lý do mua.

\begin{figure}[H]
\centering
\begin{tikzpicture}[
  ent/.style={draw, rounded corners=3pt, fill=blue!6, font=\small\sffamily, align=center, inner sep=6pt, minimum width=3cm, minimum height=1.1cm},
  rel/.style={-{Latex[length=2.5mm]}, thick, gray!60!black},
  lbl/.style={font=\scriptsize, fill=white, inner sep=1.5pt}
]
\node[ent] (cust) at (0,0) {\textbf{Khách hàng}\\ 19.119};
\node[ent] (person) at (-5,1.8) {\textbf{Hồ sơ cá nhân}\\ tuổi, thu nhập, nghề\dots};
\node[ent] (store) at (-5,-1.8) {\textbf{Hồ sơ cửa hàng}\\ loại hình, chuyên môn};
\node[ent] (terr) at (0,2.6) {\textbf{Lãnh thổ}\\ 10 vùng};
\node[ent] (soh) at (5,0) {\textbf{Đơn hàng}\\ 31.465};
\node[ent] (sod) at (5,-2.8) {\textbf{Dòng sản phẩm}\\ 121.317};
\node[ent] (prod) at (0,-2.8) {\textbf{Sản phẩm}\\ 504, 4 danh mục};
\node[ent] (reason) at (5,2.6) {\textbf{Lý do mua}\\ giá, KM, quảng cáo};
\draw[rel] (cust) -- node[lbl]{cá nhân} (person);
\draw[rel] (cust) -- node[lbl]{đại lý} (store);
\draw[rel] (cust) -- (terr);
\draw[rel] (soh) -- node[lbl]{1 khách -- n đơn} (cust);
\draw[rel] (sod) -- node[lbl]{1 đơn -- n dòng} (soh);
\draw[rel] (sod) -- (prod);
\draw[rel, dashed] (soh) -- (reason);
\end{tikzpicture}
\caption{Các thành phần trung tâm của kho dữ liệu nhìn từ góc độ khách hàng.}
\label{fig:erd}
\end{figure}

\subsection{Dữ liệu này đáng tin đến đâu?}
Toàn bộ 90 mối liên kết giữa các bảng còn nguyên vẹn, không có mã khách hay mã đơn bị trùng. 100\% khách cá nhân có đầy đủ hồ sơ khảo sát; 100\% cửa hàng có hồ sơ và có nhân viên phụ trách. Những vấn đề tìm thấy (Bảng~\ref{tab:quality}) đều nhỏ và có cách xử lý rõ ràng, nên kho dữ liệu này đủ sạch để phân tích mà không cần công đoạn làm sạch nặng.

\begin{table}[H]
\centering
\caption{Các vấn đề về chất lượng dữ liệu và cách xử lý}
\label{tab:quality}
\small
\begin{tabularx}{\textwidth}{>{\raggedright\arraybackslash}p{3.4cm} X >{\raggedright\arraybackslash}p{3.8cm}}
\toprule
\rowcolor{tblhead} \textbf{Vấn đề} & \textbf{Chi tiết} & \textbf{Cách xử lý} \\
\midrule
701 dòng \texttt{Customer} không có đơn & \texttt{PersonID} rỗng, \texttt{StoreID} có giá trị; không có đơn hàng nào & Bỏ ra khỏi tập phân tích; 19.119 khách còn lại đều có đơn \\
2 đơn có tổng tiền âm & Đơn 43678 (19 dòng hàng) và 43682 (5 dòng), ngày 31/05/2011, kênh đại lý, trạng thái 5. Toàn bộ dòng hàng của hai đơn đều có \texttt{LineTotal} âm; dữ liệu không ghi lý do (có thể là trả hàng hoặc bút toán điều chỉnh) & Giữ trong số liệu mô tả; loại khỏi tập phân tích churn vì không thể xác nhận đây là lần mua \\
315 đơn trạng thái 6 & Cột \texttt{Status} chỉ có hai giá trị: 5 (31.150 đơn) và 6 (315 đơn, tổng 1,18 triệu USD). Theo tài liệu AdventureWorks, 5 = đã giao, 6 = đã hủy & Giữ trong số liệu mô tả; không tính là lần mua trong tập phân tích churn \\
76\% đơn dồn vào 12 tháng cuối & Kênh online tăng vọt từ quý 3/2013 & Chọn mốc gắn nhãn rời bỏ sao cho khách có đủ thời gian quan sát \\
Một số trường chữ chưa chuẩn & ``Bachelors '' thừa khoảng trắng; thu nhập ghi theo khoảng chứ không phải số & Chuẩn hóa trước khi dùng làm biến phân tích \\
\bottomrule
\end{tabularx}
\end{table}

Báo cáo dùng hai tập dữ liệu và ghi rõ ở từng chỗ. \emph{Tập mô tả} là toàn bộ dữ liệu gốc, dùng cho các thống kê mô tả để giữ đúng con số của kho dữ liệu. \emph{Tập phân tích churn} loại hai đơn âm và 315 đơn hủy; Bảng~\ref{tab:before-after} cho thấy việc loại này làm mất 124 khách (những khách chỉ có đơn hủy) và khoảng 1\% đơn, dòng hàng và doanh thu.

\begin{table}[H]
\centering
\caption{Quy mô dữ liệu trước và sau khi loại đơn hủy và đơn âm (truy vấn V02e)}
\label{tab:before-after}
\footnotesize
\begin{tabularx}{\textwidth}{Xrrrrr}
\toprule
\rowcolor{tblhead} \textbf{Tập dữ liệu} & \textbf{Khách} & \textbf{Đơn} & \textbf{Dòng hàng} & \makecell[r]{\textbf{Trước thuế}\\\textbf{(USD)}} & \makecell[r]{\textbf{Tổng phải trả}\\\textbf{(USD)}} \\
\midrule
Gốc (tập mô tả) & 19.119 & 31.465 & 121.317 & 109.784.736 & 123.155.141 \\
Đã lọc (tập phân tích churn) & 18.995 & 31.148 & 120.123 & 108.748.472 & 121.990.642 \\
\bottomrule
\end{tabularx}
\end{table}
